\chapter{La construction : ordre, fraîcheur, parallélisme}

C'est le chapitre le plus important du livre, et sa dernière section est une
démonstration qu'il faut refaire soi-même au moins une fois.

\section{L'ordre}

\begin{nkretenir}
Jenga trie les projets par leurs dépendances --- un \emph{tri topologique} --- et
affiche le résultat :

\begin{nkcode}{Capture reelle}
Build Order (1 projects):
  1. Bonjour [CONSOLE_APP]
\end{nkcode}

\textbf{Cette ligne est un contrat.} Ce qui n'y figure pas ne sera pas construit,
et aucune modification de ses sources n'aura d'effet.
\end{nkretenir}

\begin{nkpiege}
\textbf{« Le build est vert » ne dit rien de ce qui a été construit.} Mesuré sur
ce dépôt : un développeur modifie un module, lance une construction, obtient
\texttt{Projects Built: 19/19} et \texttt{SUCCESS} --- et le fichier objet de sa
modification \textbf{n'existait pas}. La cible construite ne tirait pas son
module.

\emph{Un build vert prouve que ce qu'il a construit compile, pas que ce que vous
avez écrit a été construit.}

Trois vérifications, et elles sont bon marché : le \texttt{.obj} existe et est
plus récent que la source ; une chaîne caractéristique de votre modification est
dans la bibliothèque produite ; la ligne \texttt{Compiled: <votre fichier>}
apparaît dans la sortie.
\end{nkpiege}

\section{Ce qui est recompilé, et ce qui ne l'est pas}

\begin{nkretenir}
La décision se prend fichier par fichier, et voici la fonction qui la prend :

\begin{nkcode}{Jenga/Core/Builder.py:1380-1405}
if not obj.exists():
    return True

obj_mtime = obj.stat().st_mtime
if src.stat().st_mtime > obj_mtime:
    return True

dep_file = self.GetDependencyFilePath(objectFile)
if not dep_file.exists():
    return True

deps = self._ParseDependencyFile(dep_file, project)
if not deps:
    return True

for dep in deps:
    if not dep.exists():
        return True
    if dep.stat().st_mtime > obj_mtime:
        return True
\end{nkcode}

\textbf{Quatre raisons de recompiler}, dans l'ordre où elles sont testées :
l'objet n'existe pas ; la source est plus récente que lui ; le fichier de
dépendances manque ; un en-tête dont il dépend est plus récent.
\end{nkretenir}

\begin{nkretenir}
\textbf{Les en-têtes sont suivis, et c'est le compilateur qui les liste.} Le
fichier de dépendances est produit par le compilateur lui-même à la compilation
précédente : Jenga ne devine pas quels \texttt{\#include} vous avez écrits, il
lit ce que le compilateur a réellement ouvert.

C'est pourquoi la \emph{première} construction est plus lente : il n'y a pas
encore de fichier de dépendances, donc tout est recompilé.
\end{nkretenir}

\section{La décision se prend sur les DATES}

\begin{nkpiege}[Le piège central de Jenga, et il est mesurable]
Regardez à nouveau la fonction ci-dessus : elle compare des
\texttt{st\_mtime} --- des \textbf{dates}. Pas des empreintes.

Conséquence : \textbf{une source modifiée dont la date est ANCIENNE est déclarée
à jour}, et l'ancien binaire tourne.

Cela n'a rien de théorique. Une date ancienne arrive dès qu'on restaure un
fichier :

\begin{center}
\begin{tabular}{@{}ll@{}}
\toprule
\textbf{Geste} & \textbf{Ce qu'il fait à la date} \\
\midrule
\texttt{cp -p}, \texttt{rsync -t} & la \textbf{préserve} \\
extraire une archive \texttt{tar}, \texttt{zip} & la préserve \\
\texttt{Copy-Item} (PowerShell) & la préserve \\
\texttt{git checkout} & la met à \emph{maintenant} --- sûr \\
\texttt{cp} sans \texttt{-p} & maintenant --- sûr \\
\bottomrule
\end{tabular}
\end{center}
\end{nkpiege}

\subsection*{La démonstration}

Voici l'expérience, faite sur le projet du chapitre 3. À chaque étape on modifie
la source, on lance \texttt{jenga build}, et \textbf{on lance le programme} :

\begin{nkcode}{Capture reelle -- ce que le PROGRAMME affiche}
1. source = VERSION UN, construite normalement       -> VERSION UN
2. source = VERSION DEUX, date NEUVE                 -> VERSION DEUX
3. source = VERSION TROIS, date REMISE A L'ANCIENNE  -> VERSION DEUX
4. la meme source, apres `touch`                     -> VERSION TROIS
\end{nkcode}

\begin{nkretenir}
\textbf{Lisez la ligne 3 deux fois.} Le fichier sur le disque dit
\texttt{VERSION TROIS}. Le programme dit \texttt{VERSION DEUX}. La construction
a répondu \texttt{SUCCESS}.

Et la ligne 4 montre le remède : un simple \texttt{touch} --- qui ne change pas
un octet du contenu --- suffit à déclencher la recompilation.
\end{nkretenir}

\begin{nkretenir}[Le témoin choisi, et pourquoi celui-là]
\textbf{Le témoin n'est PAS la sortie de \texttt{jenga}.} Une ligne
« déjà à jour » serait déjà une réponse, et l'on ne saurait pas si elle est
juste.

Le témoin est \textbf{ce que le programme affiche} : une compilation qui n'a pas
eu lieu et une compilation sans effet se ressemblent dans un journal ; elles ne
se ressemblent pas à l'écran.
\end{nkretenir}

\begin{nkretenir}[Il existe un moteur à empreintes, et il ne sert pas]
Jenga porte un module \nkfile{Jenga/Core/Incremental.py}, 227 lignes, qui calcule
des \texttt{sha256} --- et son propre commentaire dit pourquoi :

\begin{nkcode}{Jenga/Core/Incremental.py:36-40}
def ComputeFileHash(filepath: str, algorithm: str = "sha256") -> str:
    """Calcule le hash d'un fichier. Utilise pour detecter les changements.
    Par defaut sha256, plus fiable que mtime seul."""
\end{nkcode}

\textbf{Mesuré : aucune de ses méthodes n'est appelée.} Le module est importé par
\texttt{Daemon.py} et par \texttt{Core/\_\_init\_\_.py}, et une recherche de
\texttt{Incremental.} hors du fichier lui-même rend \textbf{zéro}. La méthode
\texttt{Cache.UpdateIncremental} porte un nom voisin mais fait autre chose :
elle recharge la \emph{description}, pas les sources.

C'est un doublon dont un seul côté est vivant --- et l'on croit lire un système à
empreintes en lisant un système à dates. \emph{Une capacité écrite n'est pas une
capacité employée ; seul un appelant le prouve.}
\end{nkretenir}

\subsection*{Ce qu'il faut en faire}

\begin{center}
\begin{tabular}{@{}ll@{}}
\toprule
\textbf{Situation} & \textbf{Le geste} \\
\midrule
après toute restauration de fichiers & \texttt{touch} les fichiers restaurés \\
un doute sur la fraîcheur d'un binaire & \texttt{jenga rebuild} \\
avant une campagne de mesures & \texttt{rebuild}, puis vérifier une date \\
un correctif « sans effet » & vérifier le \texttt{.obj} avant d'accuser le code \\
\bottomrule
\end{tabular}
\end{center}

\begin{nkretenir}
\textbf{Le pire de ce piège n'est pas le temps perdu : c'est la fausse
conclusion.} On retire un correctif juste, parce qu'« il ne change rien ». Le
défaut coûte alors deux fois --- une fois à ne pas s'appliquer, une fois à faire
défaire une bonne chose.
\end{nkretenir}

\section{Le parallélisme}

\begin{nkretenir}
Les fichiers d'un même projet se compilent en parallèle ; les projets, eux,
suivent l'ordre des dépendances. C'est la seule combinaison possible : deux
projets indépendants \emph{pourraient} aller de front, mais un projet et sa
dépendance, non.

Conséquence pratique : \textbf{un graphe très en chaîne se parallélise mal.}
Trente projets dont chacun dépend du précédent se construisent en trente étapes,
quel que soit le nombre de c\oe urs.
\end{nkretenir}

\begin{nkpiege}
\textbf{En construction parallèle, l'ordre des messages n'est pas l'ordre des
événements.} Une erreur affichée après un succès peut le précéder dans le temps.

Et par défaut, \textbf{la construction s'arrête au premier échec} --- il y a donc
exactement un échec par relevé. C'est utile à savoir avant de compter les
erreurs : \texttt{Failed: 1} ne veut pas dire qu'un seul fichier est cassé, mais
qu'on s'est arrêté au premier.

\texttt{-{}-keep-going} continue et donne le compte réel.
\end{nkpiege}

\section{Nettoyer}

\begin{center}
\begin{tabular}{@{}ll@{}}
\toprule
\texttt{jenga clean} & supprime les objets et les artefacts \\
\texttt{jenga rebuild} & \texttt{clean} puis \texttt{build} \\
\bottomrule
\end{tabular}
\end{center}

\begin{nkretenir}
\textbf{\texttt{rebuild} est la réponse à toute question de fraîcheur.} Il coûte
du temps et il tranche : après lui, ce qui tourne vient de ce qui est sur le
disque.

C'est la commande à lancer avant de conclure qu'un correctif ne marche pas --- et
avant toute mesure de performance qu'on publiera.
\end{nkretenir}

\begin{nkbilan}
Jenga trie les projets par dépendances et affiche l'ordre : ce qui n'y figure pas
ne sera pas construit, et un build vert ne prouve pas que \emph{votre} fichier
l'a été. La recompilation se décide sur quatre critères, tous fondés sur des
\textbf{dates} --- pas sur des empreintes, bien qu'un moteur à \texttt{sha256}
existe dans le code, dormant. Une source restaurée avec une date ancienne est
donc déclarée à jour : c'est démontré, et le remède est \texttt{touch} ou
\texttt{rebuild}. Enfin les fichiers se parallélisent, les projets non, et la
construction s'arrête au premier échec sauf \texttt{-{}-keep-going}.
\end{nkbilan}

\begin{nkquestions}
\begin{enumerate}
\item Quelles sont les quatre raisons de recompiler un fichier ?
\item Sur quoi la décision se fonde-t-elle : dates ou empreintes ?
\item Citez trois gestes qui préservent la date d'un fichier.
\item Pourquoi \texttt{Failed: 1} ne signifie-t-il pas « un seul fichier cassé » ?
\item Que prouve un build vert, et que ne prouve-t-il pas ?
\end{enumerate}
\end{nkquestions}

\begin{nkpratique}
Reproduisez la démonstration des dates, vous-même, sur votre projet.

\begin{enumerate}
\item construisez, lancez, notez la sortie ;
\item modifiez la source normalement, reconstruisez, lancez : la sortie change ;
\item modifiez-la à nouveau \emph{en remettant l'ancienne date}
      (\texttt{os.utime} en Python, \texttt{touch -r} en shell), reconstruisez,
      lancez ;
\item \texttt{touch}, reconstruisez, lancez.
\end{enumerate}

\textbf{Notez les quatre sorties du programme}, pas celles de Jenga.
\end{nkpratique}

\begin{nkdemo}
Prouvez que le module à empreintes ne sert pas.

\begin{enumerate}
\item Cherchez dans le dépôt de Jenga les appels à \texttt{Incremental.} en
      excluant le fichier lui-même. Notez le compte.
\item Cherchez qui \emph{importe} le module. Notez que l'import existe.
\item Concluez sur la différence entre les deux comptes.
\end{enumerate}

\alerte{} \textbf{Le contrôle positif est obligatoire avant de croire un zéro} :
lancez d'abord la même recherche sur un nom dont vous savez qu'il est employé ---
\texttt{Builder.} par exemple. Si elle rend zéro aussi, c'est votre commande qui
ne sait pas chercher, et le premier zéro ne prouvait rien.
\end{nkdemo}

\begin{nklogique}
Vous corrigez un défaut, reconstruisez, et le comportement ne change pas.

\begin{enumerate}
\item Énumérez au moins cinq causes, en séparant celles où la compilation
      \emph{n'a pas eu lieu} de celles où elle a eu lieu \emph{sans effet}.
\item Rangez-les par \textbf{coût de vérification}, du moins cher au plus cher
      --- et non par probabilité.
\item Justifiez ce critère : pourquoi, quand on ne sait rien, vaut-il mieux
      ordonner par coût que par probabilité ?
\end{enumerate}
\end{nklogique}
